\documentclass{minimal}
\usepackage{xltxtra}
\usepackage{fontspec}
\usepackage{polyglossia}
\setmainlanguage{russian}
\setotherlanguage{english}
\setmainfont{Times New Roman}
\newfontfamily{\cyrillicfont}{Times New Roman}
\usepackage{amsmath}
\usepackage{tikz}
\usetikzlibrary{arrows.meta}
\begin{document}
    Если до погружения первого тела в воду на тело действовало ускорение $a$, а после погружения части первого тела в воду, система будет находиться в равносвесии, то равнодействующая сила будет равна по модулю силе Архимеда и направлена в протиположную строну. Сила Архимеда будет равна:
    $$F_{\mbox{арх}}=\rho g V$$ и направлена вверх. Равнодествующая будет такой же по модую и направлена вниз. Откуда ускорение, которое действует на систему тел: $$a=\frac{-\rho g V}{m_1 + m_2}$$
    А с другой стороны мы знаем, что это ускорение равно:
    $$F=m_1 a$$ или $$\rho g V=m_1 a$$ составим и решим уравнение:
    $$
    \begin{cases}
        a&=\frac{-\rho g V}{m_1 + m_2}\\
        \rho g V&=m_1 a
    \end{cases}
    $$
    Решая это уравнение получим:
    $$
      m_2 = 2 \cdot 0.5~\mbox{кг} = 1~\mbox{кг}
    $$
    А ускорение будет равно:
    $$
        a = \frac{1000~\frac{\mbox{кг}}{\mbox{м}^3} \cdot 9.8~\frac{\mbox{м}}{\mbox{с}^2} \cdot 1.5 \cdot 10^{-4}~\mbox{м}^2}{0.5~\mbox{кг} + 1~\mbox{кг}} = 1~\frac{\mbox{м}}{\mbox{с}^2}
    $$
\end{document}